% GENERATED FILE. Edit canonical lessons, then run npm run generate:books.
% canonical-source-hash: fnv1a64:c70028ae997f7ea7
% canonical-lessons: GE-C219-neugierig, GE-C219-nuetzlich, GE-C219-oeffentlich, GE-C219-ordentlich, GE-C219-persoenlich, GE-R219-first-pass-words-from-chapters-211-214, GE-R219-second-pass-words-from-chapters-215-219

\chapter{Curious, Useful, Public, Tidy, Personal}
\label{ch:ge-curious-useful-public-tidy-personal}
\hlchaptermodality{219}

\begin{chapteropening}
\textbf{By the end of this chapter:} \emph{I can say curious, useful, public, tidy and personal.}

Complete the last lesson of chapter 219: i can say curious, useful, public, tidy and personal.
\end{chapteropening}

\section[neugierig]{neugierig — curious}

\label{lesson:GE-C219-neugierig}



\begin{quote}
Before the new one: say the German for thin, then the German for brave.
\end{quote}

\subsection*{You'll want to know: neugierig}
\textbf{neugierig} — \textquotedblleft{}curious\textquotedblright{}.

\begin{grammarlens}[title={What you've built}]
One more word for this chapter.
\end{grammarlens}

\subsection*{Your turn}
\begin{itemize}
\raggedright
  \item \emph{Say it:} \emph{neugierig}
  \item \emph{Say it:} \emph{neugierig}, once more
  \item \emph{Say it:} say \emph{neugierig}
  \item \emph{Recall:} say the German for thin, then the German for brave, then say \emph{neugierig} again
\end{itemize}

\begin{tcolorbox}[breakable,colback=teal!4,colframe=teal!35!black,title={Before you move on}]
What does neugierig mean? (\textquotedblleft{}Curious\textquotedblright{}.) Say it once more.
\end{tcolorbox}
\section[nützlich]{nützlich — useful}

\label{lesson:GE-C219-nuetzlich}



\begin{quote}
Before the new one: say the German for brave, then the German for curious.
\end{quote}

\subsection*{You'll want to know: nützlich}
\textbf{nützlich} — \textquotedblleft{}useful\textquotedblright{}.

\begin{grammarlens}[title={What you've built}]
One more word for this chapter.
\end{grammarlens}

\subsection*{Your turn}
\begin{itemize}
\raggedright
  \item \emph{Say it:} \emph{nützlich}
  \item \emph{Say it:} \emph{nützlich}, once more
  \item \emph{Say it:} say \emph{nützlich}
  \item \emph{Recall:} say the German for brave, then the German for curious, then say \emph{nützlich} again
\end{itemize}

\begin{tcolorbox}[breakable,colback=teal!4,colframe=teal!35!black,title={Before you move on}]
What does nützlich mean? (\textquotedblleft{}Useful\textquotedblright{}.) Say it once more.
\end{tcolorbox}
\section[öffentlich]{öffentlich — public}

\label{lesson:GE-C219-oeffentlich}



\begin{quote}
Before the new one: say the German for curious, then the German for useful.
\end{quote}

\subsection*{You'll want to know: öffentlich}
\textbf{öffentlich} — \textquotedblleft{}public\textquotedblright{}.

\begin{grammarlens}[title={What you've built}]
One more word for this chapter.
\end{grammarlens}

\subsection*{Your turn}
\begin{itemize}
\raggedright
  \item \emph{Say it:} \emph{öffentlich}
  \item \emph{Say it:} \emph{öffentlich}, once more
  \item \emph{Say it:} say \emph{öffentlich}
  \item \emph{Recall:} say the German for curious, then the German for useful, then say \emph{öffentlich} again
\end{itemize}

\begin{tcolorbox}[breakable,colback=teal!4,colframe=teal!35!black,title={Before you move on}]
What does öffentlich mean? (\textquotedblleft{}Public\textquotedblright{}.) Say it once more.
\end{tcolorbox}
\section[ordentlich]{ordentlich — tidy, proper}

\label{lesson:GE-C219-ordentlich}



\begin{quote}
Before the new one: say the German for useful, then the German for public.
\end{quote}

\subsection*{You'll want to know: ordentlich}
\textbf{ordentlich} — \textquotedblleft{}tidy, proper\textquotedblright{}.

\begin{grammarlens}[title={What you've built}]
One more word for this chapter.
\end{grammarlens}

\subsection*{Your turn}
\begin{itemize}
\raggedright
  \item \emph{Say it:} \emph{ordentlich}
  \item \emph{Say it:} \emph{ordentlich}, once more
  \item \emph{Say it:} say \emph{ordentlich}
  \item \emph{Recall:} say the German for useful, then the German for public, then say \emph{ordentlich} again
\end{itemize}

\begin{tcolorbox}[breakable,colback=teal!4,colframe=teal!35!black,title={Before you move on}]
What does ordentlich mean? (\textquotedblleft{}Tidy, proper\textquotedblright{}.) Say it once more.
\end{tcolorbox}
\section[persönlich]{persönlich — personal}

\label{lesson:GE-C219-persoenlich}



\begin{quote}
Before the new one: say the German for public, then the German for tidy.
\end{quote}

\subsection*{You'll want to know: persönlich}
\textbf{persönlich} — \textquotedblleft{}personal\textquotedblright{}.

\begin{grammarlens}[title={What you've built}]
One more word for this chapter.
\end{grammarlens}

\subsection*{Your turn}
\begin{itemize}
\raggedright
  \item \emph{Say it:} \emph{persönlich}
  \item \emph{Say it:} \emph{persönlich}, once more
  \item \emph{Say it:} say \emph{persönlich}
  \item \emph{Recall:} say the German for public, then the German for tidy, then say \emph{persönlich} again
\end{itemize}

\begin{tcolorbox}[breakable,colback=teal!4,colframe=teal!35!black,title={Before you move on}]
What does persönlich mean? (\textquotedblleft{}Personal\textquotedblright{}.) Say it once more.
\end{tcolorbox}
\section[(dialogue)]{Words from chapters 211-214 — a first pass}

\label{lesson:GE-R219-first-pass-words-from-chapters-211-214}



\begin{quote}
Say the words for tidy and personal.
\end{quote}

\subsection*{Your turn}
Say these aloud:

\begin{itemize}
\raggedright
  \item anxious, tiring, excited, excellent, working
  \item cloudy, blond, honest, urgent, lonely
  \item jealous, successful, having a cold, fantastic, damp
  \item cheerful, terrible, patient, stingy, cosy
\end{itemize}

\begin{tcolorbox}[breakable,colback=teal!4,colframe=teal!35!black,title={Before you move on}]
Anxious? (\textbf{ängstlich}.) Tiring? (\textbf{anstrengend}.) Excited? (\textbf{aufgeregt}.) Excellent? (\textbf{ausgezeichnet}.) Working? (\textbf{berufstätig}.) Cloudy? (\textbf{bewölkt}.) Blond? (\textbf{blond}.) Honest? (\textbf{ehrlich}.) Urgent? (\textbf{eilig}.) Lonely? (\textbf{einsam}.) Jealous? (\textbf{eifersüchtig}.) Successful? (\textbf{erfolgreich}.) Having a cold? (\textbf{erkältet}.) Fantastic? (\textbf{fantastisch}.) Damp? (\textbf{feucht}.) Cheerful? (\textbf{fröhlich}.) Terrible? (\textbf{furchtbar}.) Patient? (\textbf{geduldig}.) Stingy? (\textbf{geizig}.) Cosy? (\textbf{gemütlich}.)
\end{tcolorbox}
\section[(dialogue)]{Words from chapters 215-219 — a second pass}

\label{lesson:GE-R219-second-pass-words-from-chapters-215-219}



\begin{quote}
Say the words for fair and divorced.
\end{quote}

\subsection*{Your turn}
Say these aloud:

\begin{itemize}
\raggedright
  \item fair, divorced, at the same time, generous, good value
  \item thorough, frequent, wonderful, helpful, intelligent
  \item broken, clever, free of charge, strong, single
  \item tasty, kind, relaxed, thin, brave
  \item curious, useful, public, tidy, personal
\end{itemize}

\begin{tcolorbox}[breakable,colback=teal!4,colframe=teal!35!black,title={Before you move on}]
Fair? (\textbf{gerecht}.) Divorced? (\textbf{geschieden}.) At the same time? (\textbf{gleichzeitig}.) Generous? (\textbf{großzügig}.) Good value? (\textbf{günstig}.) Thorough? (\textbf{gründlich}.) Frequent? (\textbf{häufig}.) Wonderful? (\textbf{herrlich}.) Helpful? (\textbf{hilfsbereit}.) Intelligent? (\textbf{intelligent}.) Broken? (\textbf{kaputt}.) Clever? (\textbf{klug}.) Free of charge? (\textbf{kostenlos}.) Strong? (\textbf{kräftig}.) Single? (\textbf{ledig}.) Tasty? (\textbf{lecker}.) Kind? (\textbf{lieb}.) Relaxed? (\textbf{locker}.) Thin? (\textbf{mager}.) Brave? (\textbf{mutig}.) Curious? (\textbf{neugierig}.) Useful? (\textbf{nützlich}.) Public? (\textbf{öffentlich}.) Tidy? (\textbf{ordentlich}.) Personal? (\textbf{persönlich}.)
\end{tcolorbox}
